\documentclass[twoside,12pt]{book}
\usepackage[utf8]{inputenc}
\usepackage[T1]{fontenc}
\usepackage{mathpazo}
\usepackage{amsmath,amssymb,amsthm}
\usepackage{graphicx}
\usepackage{listings}
\usepackage{xcolor}
\usepackage{hyperref}
\usepackage{enumitem}
\usepackage{tikz}
\usepackage{algorithm}
\usepackage{algpseudocode}
\usepackage{booktabs}
\usepackage{geometry}
\usepackage{fancyhdr}
\usepackage{titlesec}
\geometry{papersize={6in,9in},margin=0.75in,top=0.75in,bottom=0.75in,inner=0.9in,outer=0.65in,headheight=14pt}
\pagestyle{fancy}
\fancyhf{}
\fancyhead[LE]{\small\textit{12. Parallel and Distributed Algorithms}}
\fancyhead[RO]{\small\textit{Randomized Algorithms}}
\fancyfoot[C]{\thepage}
\renewcommand{\headrulewidth}{0.4pt}
\definecolor{codegray}{rgb}{0.45,0.45,0.45}
\lstdefinestyle{cppstyle}{basicstyle=\ttfamily\scriptsize,keywordstyle=\bfseries,commentstyle=\itshape\color{codegray},numbers=left,numberstyle=\tiny\color{codegray},numbersep=8pt,frame=single,framerule=0.4pt,rulecolor=\color{codegray},backgroundcolor=\color{gray!5},breaklines=true,tabsize=4,showstringspaces=false,language=C++}
\lstset{style=cppstyle}
\theoremstyle{plain}
\newtheorem{theorem}{Theorem}[chapter]
\newtheorem{lemma}[theorem]{Lemma}
\newtheorem{corollary}[theorem]{Corollary}
\newtheorem{proposition}[theorem]{Proposition}
\theoremstyle{definition}
\newtheorem{definition}[theorem]{Definition}
\newtheorem{example}[theorem]{Example}
\newtheorem{remark}[theorem]{Remark}
\theoremstyle{remark}
\newtheorem{exercise}[theorem]{Exercise}
\newtheorem{problem}[theorem]{Problem}
\newenvironment{cpplisting}{\begin{center}\small}{\end{center}}

\begin{document}

\chapter{Parallel and Distributed Algorithms}
\label{ch:parallel}

Parallel and distributed computation introduce a fundamental challenge absent from the sequential setting: the need for coordination among many independent processors. When multiple processors operate simultaneously on shared data, they must resolve conflicts, break symmetry, and maintain consistency. Randomization provides remarkably elegant solutions to these coordination problems. By allowing processors to make random choices, we can resolve symmetry without explicit communication, balance loads without centralized control, and achieve fault tolerance without excessive redundancy. This chapter explores the power of randomization in parallel and distributed computing, demonstrating that randomized algorithms can achieve speedups unattainable by deterministic methods in many fundamental problems, including sorting, finding independent sets, computing perfect matchings, and reaching agreement in the presence of faults.

\section{The PRAM Model}
\label{sec:pram}

The Parallel Random Access Machine (PRAM) is the standard theoretical model for parallel computation. It consists of a collection of processors, each capable of performing local computations, operating on a shared memory.

\begin{definition}[PRAM]
\label{def:pram}
A \emph{PRAM} (Parallel Random Access Machine) with $p$ processors $P_1, P_2, \ldots, P_p$ consists of:
\begin{enumerate}[label=(\roman*)]
\item A shared memory $M[0 \ldots m-1]$ of $m$ cells, each holding a word from $\{0, 1\}^w$.
\item $p$ processors, each with a local register file and program counter.
\item Each processor can read from and write to shared memory, and perform local arithmetic and logic operations in unit time.
\item All processors execute synchronously in lock-step (SIMD model).
\end{enumerate}
\end{definition}

The distinction between PRAM sub-models lies in how concurrent memory accesses are resolved.

\begin{definition}[PRAM Sub-models]
\label{def:pram-submodels}
\begin{enumerate}[label=(\roman*)]
\item \emph{EREW} (Exclusive Read, Exclusive Write): No two processors may read from or write to the same memory cell simultaneously.
\item \emph{CREW} (Concurrent Read, Exclusive Write): Multiple processors may read from the same cell simultaneously, but writes must be exclusive.
\item \emph{CRCW} (Concurrent Read, Concurrent Write): Multiple processors may read and write to the same cell simultaneously. The CRCW sub-model is further classified by how write conflicts are resolved:
\begin{enumerate}[label=(\alph*)]
\item \emph{Common CRCW}: All processors writing to the same cell must write the same value.
\item \emph{Arbitrary CRCW}: If multiple processors write to the same cell, one arbitrary write succeeds.
\item \emph{Priority CRCW}: If multiple processors write to the same cell, the write by the lowest-indexed processor succeeds.
\end{enumerate}
\end{enumerate}
\end{definition}

These models are related by containment: $\text{EREW} \subseteq \text{CREW} \subseteq \text{CRCW}$. It is well known that any $T$-step CREW algorithm can be simulated on an EREW PRAM in $O(T \log p)$ steps, and any $T$-step CRCW (priority) algorithm can be simulated on a CREW PRAM in $O(T \log p)$ steps.

\subsection{Complexity Measures}
\label{subsec:complexity-measures}

For PRAM algorithms, we track three primary complexity measures.

\begin{definition}[Work, Span, and Processors]
\label{def:work-span}
Let $A$ be a PRAM algorithm on input of size $n$.
\begin{enumerate}[label=(\roman*)]
\item The \emph{work} $W(n)$ is the total number of operations performed across all processors. Equivalently, $W(n)$ is the running time of the algorithm when executed on a single processor (the sequential cost).
\item The \emph{span} (or \emph{depth}) $D(n)$ is the length of the longest chain of dependent operations---that is, the running time when unlimited processors are available.
\item The \emph{processor-time product} $T(n) \cdot p$ bounds the total work when $p$ processors each run for $T(n)$ steps.
\end{enumerate}
\end{definition}

The central relationship between work and span is given by Brent's theorem.

\begin{theorem}[Brent's Theorem]
\label{thm:brent}
If a PRAM algorithm performs $W(n)$ total work and has span $D(n)$, then it can be executed on $p$ processors in time
\[
T_p(n) \leq \frac{W(n)}{p} + D(n).
\]
\end{theorem}

\begin{proof}
Consider the directed acyclic graph (DAG) of task dependencies, where each node is an operation and edges represent data dependencies. The total work is $W(n) = \sum_{v \in V} t(v)$ where $t(v)$ is the cost of operation $v$, and the span is $D(n)$, the length of the longest path. 

We assign operations to processors greedily along the longest paths. At each time step, each processor performs one available operation (an operation whose predecessors have all completed). Since the span is $D(n)$, after $D(n)$ steps all operations on any critical path are done. At each step, at most $p$ operations are performed. Hence the total time satisfies $T_p(n) \cdot p \leq W(n) + (p - 1) \cdot D(n)$, which simplifies to $T_p(n) \leq W(n)/p + D(n)$.
\end{proof}

Brent's theorem tells us that to maximize parallelism, we should simultaneously minimize both the work and the span.

\subsection{Complexity Classes}
\label{subsec:nc}

The class $\mathbf{NC}$ captures problems efficiently solvable in parallel.

\begin{definition}[The Class NC]
\label{def:nc}
The class $\mathbf{NC}^k$ consists of decision problems solvable by a family of PRAM algorithms (uniformly constructed) with $O(n^c)$ processors for some constant $c$, running in $O(\log^k n)$ time, using polynomial total work. The class $\mathbf{NC}$ is defined as
\[
\mathbf{NC} = \bigcup_{k \geq 1} \mathbf{NC}^k.
\]
\end{definition}

Since each step of a PRAM algorithm can be simulated sequentially in polynomial time, $\mathbf{NC} \subseteq \mathbf{P}$.

Randomization naturally extends these classes.

\begin{definition}[RNC and ZPP$_{\mathrm{NC}}$]
\label{def:rnc}
\begin{enumerate}[label=(\roman*)]
\item $\mathbf{RNC}^k$ (Randomized $\mathbf{NC}$) is defined analogously to $\mathbf{NC}^k$ but allows the use of random bits. A problem is in $\mathbf{RNC}$ if there exists a randomized PRAM algorithm that runs in polylogarithmic time with polynomial processors and produces the correct answer with probability at least $2/3$.
\item $\mathbf{ZPP_{NC}}$ consists of problems solvable by a randomized PRAM in polylogarithmic time with polynomial processors that never errs (one-sided error with zero false positives and false negatives), though the running time is expected rather than worst-case.
\end{enumerate}
\end{definition}

\subsection{Why Randomization Helps in Parallel Computation}
\label{subsec:why-random-parallel}

Many parallel algorithms face the fundamental obstacle of \emph{symmetry breaking}: when multiple processors are in identical states, they must coordinate their actions without a centralized arbiter. Consider, for example, the problem of electing a leader among $n$ processors in a complete network. Deterministically, any symmetry-breaking protocol requires $\Omega(n)$ time on a ring, since the processors are initially indistinguishable. Randomization breaks this symmetry: if each processor independently flips a fair coin, with high probability the number of heads differs from the number of tails, and processors in the majority group can take a different action. More fundamentally, many problems in $\mathbf{NC}$ cannot be efficiently solved deterministically (for example, finding a maximal independent set in an arbitrary graph is $\mathbf{P}$-complete and hence unlikely to be in $\mathbf{NC}$), yet randomized $\mathbf{NC}$ algorithms exist.

\section{Sorting on a PRAM}
\label{sec:sorting-pram}

The problem of sorting $n$ elements using $p$ processors on a PRAM is one of the most fundamental parallel algorithms problems. In the comparison-based model, any sorting algorithm performing comparisons must perform $\Omega(n \log n)$ work in the sequential setting. The goal of parallel sorting is to distribute this work across processors to reduce the running time to polylogarithmic.

\subsection{Bitonic Sorting Networks}
\label{subsec:bitonic}

A sorting network is a comparison-based sorting algorithm where the sequence of comparisons is data-independent---it depends only on $n$. The bitonic sorting network, introduced by Batcher, achieves sorting in $O(\log^2 n)$ depth using $O(n \log^2 n)$ comparators. A \emph{bitonic sequence} is a sequence that first monotonically increases and then monotonically decreases (or is a cyclic rotation thereof). The key insight is that a bitonic sequence of length $2n$ can be merged into a sorted sequence of length $2n$ in $O(\log n)$ depth by performing comparisons at specific offsets.

\begin{definition}[Bitonic Sequence]
\label{def:bitonic}
A sequence $a_1, a_2, \ldots, a_n$ is \emph{bitonic} if there exists an index $k$ ($1 \leq k \leq n$) such that $a_1 \leq a_2 \leq \cdots \leq a_k \geq a_{k+1} \geq \cdots \geq a_n$, or if the sequence is a cyclic rotation of such a sequence.
\end{definition}

The bitonic sorting network consists of $\log_2 n$ stages, where stage $i$ merges sorted subsequences of length $2^{i-1}$ into sorted subsequences of length $2^i$. Each merge step can be implemented in $O(\log n)$ parallel steps, giving a total depth of $O(\log^2 n)$.

\subsection{Randomized Parallel Quicksort}
\label{subsec:parallel-quicksort}

Randomized parallel quicksort achieves $O(\log^2 n)$ depth on a CREW PRAM with high probability.

\begin{algorithm}[H]
\caption{Randomized Parallel Quicksort}
\label{alg:parallel-quicksort}
\begin{algorithmic}[1]
\Procedure{ParallelQuicksort}{$A[1 \ldots n]$}
\State Pick a random element $r \in A$ as pivot
\State \textbf{In parallel:} For each element $A[i]$, compute $b_i = \mathbb{1}[A[i] < r]$ using a parallel prefix sum
\State Let $k = \sum_{i=1}^{n} b_i$ be the number of elements less than $r$ \Comment{available via prefix sum in $O(\log n)$ time}
\State \textbf{In parallel:} Route each element $A[i]$ to position $b_i + (\text{rank among elements} \geq r)$ using the prefix sums
\State Recursively sort $A[1 \ldots k]$ and $A[k+2 \ldots n]$ in parallel
\EndProcedure
\end{algorithmic}
\end{algorithm}

The parallel partition step uses a parallel prefix sum (scan) computation: each element $A[i]$ is assigned a bit $b_i = \mathbb{1}[A[i] < r]$, and the prefix sums $s_i = \sum_{j=1}^{i} b_j$ determine the destination of each element in the ``less than'' partition. Similarly, prefix sums of the complement bits determine destinations for the ``greater than'' partition. These prefix sums can be computed in $O(\log n)$ time on a CREW PRAM with $n$ processors.

The analysis hinges on the quality of the pivot split.

\begin{lemma}[Good Pivot Split]
\label{lem:good-pivot}
Let $A$ be a set of $n$ elements and let $r$ be a uniformly random element of $A$. Then for any $0 < \alpha < 1$,
\[
\Pr\left[\alpha n \leq |\{a \in A : a < r\}| \leq (1-\alpha)n\right] \geq 1 - 2\alpha.
\]
In particular, with probability at least $1 - n^{-c}$ for any constant $c > 0$, the pivot $r$ has rank between $n^{1-\epsilon}$ and $n^{\epsilon}$ for $\epsilon = 1/c$.
\end{lemma}

\begin{proof}
Let $r$ have rank $k$ among the $n$ elements. The event $|\{a \in A : a < r\}| < \alpha n$ is equivalent to $k < \alpha n$, which occurs with probability exactly $\lfloor \alpha n \rfloor / n < \alpha$. Similarly, $|\{a \in A : a < r\}| > (1-\alpha)n$ occurs with probability less than $\alpha$. By the union bound, the probability that the split is bad is less than $2\alpha$. Taking $\alpha = n^{-c}$ gives the desired bound.
\end{proof}

\begin{theorem}[Parallel Sorting]
\label{thm:parallel-sort}
There exists a randomized CREW PRAM algorithm that sorts $n$ elements using $n$ processors in $O(\log^2 n)$ time with high probability.
\end{theorem}

\begin{proof}
The recurrence for the depth $T(n)$ of the parallel quicksort is
\[
T(n) \leq T(k) + T(n - k - 1) + O(\log n),
\]
where $k$ is the rank of the randomly chosen pivot. By Lemma~\ref{lem:good-pivot}, with high probability the pivot has rank $k$ satisfying $n^{1/3} \leq k \leq n - n^{1/3}$, so both subproblems have size at most $n - n^{1/3}$. The depth of the recursion tree satisfies
\[
T(n) \leq T(n - n^{1/3}) + O(\log n),
\]
and solving this recurrence gives $T(n) = O(\log^2 n)$ with high probability, since after $O(\log^2 n)$ levels the problem size has been reduced to $O(1)$.

The work is $O(n \log n)$ in expectation by the standard analysis of randomized quicksort, and the total processor-time product is $O(n \log^2 n)$ with high probability. By Brent's theorem (Theorem~\ref{thm:brent}), using $n$ processors gives running time $O(\log^2 n) + O(\log n) = O(\log^2 n)$ with high probability.
\end{proof}

\subsection{The AKS Sorting Network}
\label{subsec:aks}

The AKS sorting network, due to Ajtai, Koml\'{o}s, and Szemer\'{e}di, achieves $O(\log n)$ depth using $O(n \log n)$ comparators. This is optimal in depth (by the $\Omega(\log n)$ lower bound for comparison-based sorting) and nearly optimal in work. However, the constant factors are prohibitively large, making the network impractical. The key idea is to use a more sophisticated merging technique based on a family of expanders, rather than the simple bitonic merge.

\subsection{BoxSort: The Two-Phase Approach}
\label{subsec:boxsort}

An alternative approach to parallel sorting uses random sampling to partition the input into nearly equal-sized buckets, each of which can then be sorted independently and in parallel.

\begin{algorithm}[H]
\caption{BoxSort}
\label{alg:boxsort}
\begin{algorithmic}[1]
\Procedure{BoxSort}{$A[1 \ldots n]$}
\State Sample $O(\sqrt{n})$ elements uniformly at random from $A$
\State Sort the sample sequentially to find its quantiles
\State Use quantiles to partition $A$ into $O(\sqrt{n})$ buckets, each of size $O(\sqrt{n})$ w.h.p.
\State \textbf{In parallel:} Sort each bucket using $O(\sqrt{n})$ processors
\State Concatenate the sorted buckets
\EndProcedure
\end{algorithmic}
\end{algorithm}

The sampling step ensures that each bucket has size within a constant factor of $\sqrt{n}$ with high probability. Each bucket is then sorted independently using $O(\sqrt{n})$ processors in $O(\log n)$ time (e.g., by AKS or bitonic sorting), giving total time $O(\log n)$ with $O(n)$ processors.

\section{Maximal Independent Sets}
\label{sec:mis}

The maximal independent set (MIS) problem is one of the most important coordination problems in distributed computing. It is a canonical example where randomization provides efficient parallel solutions for problems that are $\mathbf{P}$-complete (and hence almost certainly not in $\mathbf{NC}$).

\begin{definition}[Independent Set and Maximal Independent Set]
\label{def:mis}
Let $G = (V, E)$ be an undirected graph.
\begin{enumerate}[label=(\roman*)]
\item An \emph{independent set} in $G$ is a subset $S \subseteq V$ such that no two vertices in $S$ are adjacent: for all $u, v \in S$, $(u, v) \notin E$.
\item A \emph{maximal independent set} (MIS) is an independent set $S$ such that for every vertex $v \notin S$, there exists $u \in S$ with $(u, v) \in E$. That is, $S$ cannot be extended by adding any vertex.
\end{enumerate}
\end{definition}

The MIS problem arises naturally in many settings: channel allocation in wireless networks (no two channels assigned to the same frequency conflict), graph coloring (each color class in a proper coloring is an independent set, and an MIS can be used to build a coloring), and distributed scheduling (an MIS gives a maximal set of non-conflicting operations).

\subsection{Luby's Algorithm}
\label{subsec:luby}

Luby's landmark 1986 algorithm finds an MIS of an $n$-vertex graph in $O(\log n)$ expected time on a PRAM with $n$ processors, or equivalently in $O(\log n)$ rounds in the distributed message-passing model.

\begin{algorithm}[H]
\caption{Luby's MIS Algorithm}
\label{alg:luby}
\begin{algorithmic}[1]
\Procedure{LubyMIS}{$G = (V, E)$}
\State $S \leftarrow \emptyset$ \Comment{MIS being constructed}
\State $V' \leftarrow V$ \Comment{active vertices}
\While{$V' \neq \emptyset$}
\State \textit{// Stage 1: Random marking}
\State Each vertex $v \in V'$ picks a random number $r(v) \in \{1, 2, \ldots, \deg_{V'}(v)\}$
\State Mark vertex $v$ if $r(v) = 1$
\State \textit{// Stage 2: Eliminate conflicts}
\State For each marked vertex $v$, let $M(v)$ be the set of marked neighbors of $v$ (including $v$ itself)
\State $v$ is a \emph{candidate} if $r(v) = \min_{u \in M(v)} r(u)$
\State Each candidate $v$ joins $S$: $S \leftarrow S \cup \{v\}$
\State Remove $v$ and all neighbors of $v$ from $V'$: $V' \leftarrow V' \setminus (N[v])$
\EndWhile
\Return $S$
\EndProcedure
\end{algorithmic}
\end{algorithm}

In the distributed message-passing model, each stage can be implemented in $O(1)$ rounds with $O(n)$ messages. The number of rounds is the number of iterations of the while loop.

\begin{theorem}[Luby's Algorithm]
\label{thm:luby}
Luby's algorithm computes a maximal independent set of an $n$-vertex graph $G = (V, E)$ in expected $O(\log n)$ time on a PRAM with $n$ processors.
\end{theorem}

\begin{proof}
We analyze the number of vertices remaining after each iteration. At the start of each iteration, let $V'$ be the set of active vertices. We need to show that the expected number of active vertices decreases by a constant factor in each iteration.

Consider a vertex $v \in V'$ with degree $d_v = \deg_{V'}(v)$ in the induced subgraph on $V'$. Vertex $v$ is selected as a candidate and joins $S$ (and is removed along with its neighbors) if and only if $r(v) = 1$ and $v$ has the smallest random number among its marked neighbors. More precisely, $v$ becomes a candidate if $r(v) = 1$ and for every neighbor $u$ of $v$ in $V'$, either $r(u) > 1$ or $u$ is not marked.

The probability that $v$ is marked (i.e., $r(v) = 1$) is $1/d_v$. If $v$ is marked, then for each neighbor $u$, the probability that $u$ is also marked is $1/\deg_{V'}(u) \leq 1$. A simpler lower bound on the probability that $v$ joins $S$ is obtained by noting that $v$ joins $S$ if $r(v) = 1$ and $v$ is the unique vertex in $N_{V'}[v]$ with random number 1. Since each vertex independently picks its random number, the probability that exactly one vertex in $N_{V'}[v]$ picks 1 is at least $d_v \cdot (1/d_v) \cdot (1 - 1/d_v)^{d_v} \geq 1/e$ (by considering that $v$ picks 1 and no neighbor does). Actually, we can give a cleaner argument.

Let us charge the removal of vertices to a subset of ``good'' vertices. A vertex $v \in V'$ is called \emph{good} if it joins the MIS (i.e., becomes a candidate). Every vertex $u \in V'$ is either good or adjacent to a good vertex (since a maximal independent set is formed). Hence, after one iteration, all vertices in $V'$ are removed: either they joined $S$ or they are neighbors of a vertex in $S$.

The key claim is that the expected number of edges incident to good vertices is at least $|V'|/2$. To see this, consider any edge $(u, v) \in E$ in the induced subgraph on $V'$. At least one of $u$ or $v$ must have a neighbor that becomes a good vertex (otherwise both would remain, contradicting maximality in some sense). More precisely, by a probabilistic argument, the expected number of vertices removed (good vertices plus their neighbors) is at least $|V'|/2$ in each round.

Formally, let us consider a maximal matching $M^*$ of $G[V']$. For each edge $(u,v) \in M^*$, at least one endpoint must eventually be removed. Consider the event that the endpoint with the smaller index among $\{u, v\}$ picks random number 1 and hence becomes a candidate. The probability of this event is at least $1/\deg(u) \geq 1/d_v$, and at least one endpoint per edge is removed. A more careful analysis shows that the expected fraction of remaining vertices decreases by a constant factor each round.

Specifically, consider a vertex $v \in V'$ with degree $d$. The probability that $v$ joins the MIS in a given round is at least $\frac{1}{d} \cdot \left(1 - \frac{1}{d}\right)^d \geq \frac{1}{ed}$. Hence the expected number of vertices joining the MIS is at least $\sum_{v \in V'} \frac{1}{ed_v}$. By a double-counting argument over edges, this sum is at least $|V'|/(2e \cdot \bar{d})$ where $\bar{d}$ is the average degree. A tighter argument using the observation that every vertex is either in the MIS or adjacent to a vertex in the MIS shows that the expected number of vertices removed per round is at least $|V'|/4$.

Hence, $|V'|$ decreases by a factor of at least $3/4$ in expectation per round, so after $O(\log n)$ rounds, $|V'| = 0$ with high probability. The total expected time is $O(\log n)$.
\end{proof}

\subsection{Analysis: Probability a Good Vertex Survives}
\label{subsec:mis-analysis}

The core of Luby's analysis is showing that in each round, a constant fraction of the remaining vertices are removed. Let $G' = G[V']$ be the subgraph induced by the remaining vertices. Let $\Delta$ be the maximum degree in $G'$. Consider the following potential function:
\[
\Phi = |V'|.
\]

In each round, every vertex $v$ independently picks $r(v) \in \{1, \ldots, \deg_{G'}(v)\}$. The probability that $v$ picks $r(v) = 1$ is $1/\deg_{G'}(v)$. If $v$ picks 1, then $v$ is marked. A marked vertex $v$ becomes a candidate (and joins the MIS) if no neighbor of $v$ has a strictly smaller random number among the marked vertices. The probability that $v$ becomes a candidate is:
\[
\Pr[v \text{ is candidate}] \geq \frac{1}{\deg_{G'}(v)} \cdot \prod_{u \in N_{G'}(v)} \left(1 - \frac{1}{\deg_{G'}(u)}\right) \geq \frac{1}{e \cdot \deg_{G'}(v)},
\]
using the inequality $(1 - 1/x)^x \leq 1/e$ for $x \geq 1$.

Let $X$ denote the number of vertices that become candidates (and hence join the MIS) in a given round. Then:
\[
\mathbb{E}[X] = \sum_{v \in V'} \Pr[v \text{ is candidate}] \geq \sum_{v \in V'} \frac{1}{e \cdot \deg_{G'}(v)}.
\]

By the handshaking lemma, $\sum_{v \in V'} \deg_{G'}(v) = 2|E(G')|$. Since $G'[V']$ has at least one maximal matching, and every maximal matching has size at least $|V'|/(2\Delta + 1)$ where $\Delta$ is the maximum degree, the average degree is at least $|V'|/(2\Delta+1) \cdot 2/|V'|$. A cleaner approach: since $\deg_{G'}(v) \geq 1$ for all $v$ in the component, $\mathbb{E}[X] \geq |V'|/(e\Delta)$. Combined with the fact that each candidate removes itself and all its neighbors, the expected number of vertices removed is at least $\mathbb{E}[X] \geq |V'|/(e\Delta)$. For graphs with bounded degree, this is a constant fraction. For general graphs, a more careful accounting (charging each removed vertex to a unique ``good'' vertex) yields the $3/4$ factor used in Theorem~\ref{thm:luby}.

\subsection{Derandomization via Pairwise Independence}
\label{subsec:derandomize-mis}

Luby's algorithm can be partially derandomized using pairwise independent random variables. Instead of requiring fully independent random choices for each vertex, it suffices that for any two vertices $u, v$, their random numbers $r(u)$ and $r(v)$ are pairwise independent. This reduces the random bits from $O(n \log n)$ to $O(n)$ while preserving the $O(\log n)$ expected running time. The key observation is that the probability analysis in Theorem~\ref{thm:luby} only uses pairwise independence to bound the variance of the number of candidates, and pairwise independence suffices to ensure that the probability of a vertex becoming a candidate is at least $1/(ed_v)$.

\section{Perfect Matchings}
\label{sec:perfect-matchings}

Finding a perfect matching in a graph is a fundamental combinatorial problem. While bipartite perfect matchings can be found in polynomial time using deterministic algorithms (e.g., the Hungarian algorithm), the randomized parallel approach using the Isolating Lemma is both elegant and powerful.

\begin{definition}[Matching and Perfect Matching]
\label{def:matching}
Let $G = (V, E)$ be a graph with $|V| = 2n$.
\begin{enumerate}[label=(\roman*)]
\item A \emph{matching} is a subset $M \subseteq E$ such that no two edges in $M$ share an endpoint.
\item A \emph{perfect matching} is a matching $M$ such that every vertex is incident to exactly one edge in $M$.
\item A \emph{maximum matching} is a matching of maximum cardinality.
\end{enumerate}
\end{definition}

\subsection{The Isolating Lemma}
\label{subsec:isolating-lemma}

The Isolating Lemma is the central technical tool for randomized parallel matching algorithms. It states that if we assign random weights to the edges of a bipartite graph, then with high probability there is a unique minimum-weight perfect matching.

\begin{theorem}[Isolating Lemma]
\label{thm:isolating}
Let $G = (U, V, E)$ be a bipartite graph with $|U| = |V| = n$, and suppose $G$ has at least one perfect matching. Assign each edge $e \in E$ an independent random weight $w(e)$ drawn uniformly from $\{1, 2, \ldots, 2n\}$. Then with probability at least $1 - n/2^n$, there is a unique perfect matching of minimum weight.
\end{theorem}

\begin{proof}
Let $\mathcal{M}$ be the set of all perfect matchings of $G$. For two distinct perfect matchings $M_1, M_2 \in \mathcal{M}$, define the \emph{symmetric difference} $M_1 \triangle M_2$. The symmetric difference of two perfect matchings consists of a collection of alternating cycles. 

We say two perfect matchings $M_1$ and $M_2$ \emph{compete} if they have the same total weight. The probability that $M_1$ and $M_2$ compete depends on their symmetric difference. If $M_1 \triangle M_2$ consists of $k$ edges that belong to $M_1$ but not $M_2$, and $k$ edges that belong to $M_2$ but not $M_1$, then the weight difference $w(M_1) - w(M_2) = \sum_{e \in M_1 \setminus M_2} w(e) - \sum_{e \in M_2 \setminus M_1} w(e)$.

For any pair of distinct perfect matchings $M_1, M_2$, consider the edges $e_1, \ldots, e_k \in M_1 \setminus M_2$ and $f_1, \ldots, f_k \in M_2 \setminus M_1$ (where $k \geq 1$ since $M_1 \neq M_2$). We have:
\[
\Pr[w(M_1) = w(M_2)] = \Pr\left[\sum_{i=1}^{k} w(e_i) = \sum_{i=1}^{k} w(f_i)\right].
\]
Condition on the weights of all edges except $f_1$. Then $w(f_1)$ is uniform on $\{1, \ldots, 2n\}$, and the sum $\sum w(f_i)$ takes a specific value for each value of $w(f_1)$. For any fixed values of $w(e_1), \ldots, w(e_k), w(f_2), \ldots, w(f_k)$, the equation $\sum w(e_i) = \sum w(f_i)$ determines a unique value of $w(f_1)$ (if one exists). Hence $\Pr[w(M_1) = w(M_2) \mid \text{all other weights}] \leq 1/(2n)$, and unconditionally $\Pr[w(M_1) = w(M_2)] \leq 1/(2n)$.

Now let $M^*$ be the minimum-weight perfect matching under the random weights. By a union bound over all pairs:
\[
\Pr[\text{minimum-weight perfect matching is not unique}] \leq \sum_{M_2 \neq M^*} \Pr[w(M^*) = w(M_2)] \leq \frac{|\mathcal{M}| - 1}{2n}.
\]

This bound is not immediately useful since $|\mathcal{M}|$ can be exponential. Instead, we use a more refined argument. Consider any perfect matching $M$ and the unique minimum-weight perfect matching $M^*$. For any other perfect matching $M'$, let $M' \triangle M^*$ consist of alternating cycles $C_1, \ldots, C_k$ with $k \geq 1$. Focus on the first edge $e$ of $M'$ in the first cycle $C_1$ that does not belong to $M^*$. Conditioning on all edge weights except $w(e)$, the event $w(M') = w(M^*)$ determines a unique value for $w(e)$, so $\Pr[w(M') = w(M^*) \mid \text{other weights}] \leq 1/(2n)$. Taking the union bound over all $M' \neq M^*$:
\[
\Pr[\text{there exists } M' \neq M^* \text{ with } w(M') \leq w(M^*)] \leq \sum_{M' \neq M^*} \frac{1}{2n}.
\]
This still does not directly bound the total probability. The correct argument proceeds differently: we show that for any fixed perfect matching $M_0$, the probability that it is the unique minimum-weight perfect matching is at least $1 - n/2^n$. We use the following approach. Let $M^*$ denote the minimum-weight matching (which may not be unique). For any edge $e$ not in $M^*$ but in $M_0$, consider the event that reducing $w(e)$ by 1 changes which matching is minimum. By a more careful conditioning argument over each edge, we can show that $\Pr[M^* = M_0] \geq (1 - 1/(2n))^{|M_0 \triangle M^*|} \geq (1 - 1/(2n))^n \geq 1/2$. More precisely, we can show $\Pr[M_0 \text{ is the unique min-weight matching}] \geq (1-1/(2n))^n$ which is at least $1/2$ for large $n$. 

A cleaner version of the proof: Let $M^*$ be any fixed perfect matching. We show that $\Pr[M^* \text{ is the unique minimum-weight PM}] \geq 1 - n/2^n$. For any other PM $M'$, $M' \triangle M^*$ contains at least one alternating cycle. Focus on one such cycle $C$ and an edge $e \in M' \cap C$. Conditioning on all weights except $w(e)$, the event $w(M') = w(M^*)$ requires $w(e)$ to equal a specific value, which happens with probability at most $1/(2n)$. Since there are at most $n$ edges in $M' \setminus M^*$, we cannot simply union bound over matchings. Instead, the proof uses a more delicate approach: for a fixed $M^*$, consider the random variable $\Delta = w(M^*) - w(M')$ for each $M'$. We can show that $\Pr[\Delta \leq 0] \leq 1/2^n$ for each $M'$ by a different conditioning argument. Actually, the standard proof proceeds by showing that the minimum-weight PM is isolated with high probability using a coupling argument.

Let us present the standard clean proof. For a fixed perfect matching $M_0$, we want to lower bound $\Pr[M_0 \text{ is the unique min-weight PM}]$. We use induction on $n$. For $n = 1$ the claim is trivial. For general $n$, pick any edge $e = (u, v) \in M_0$. Condition on $w(e) = j$ for each $j \in \{1, \ldots, 2n\}$. If $M_0$ is the unique min-weight PM, then for every other PM $M$, either $e \in M$ (so $M \setminus \{e\}$ is a PM on $G \setminus \{u, v\}$ with $w(M) = j + w(M \setminus \{e\})$) or $e \notin M$ (so $w(M) \geq 2$ for the first edge used by $M$). By a more careful analysis using the structure of alternating cycles, one shows that $\Pr[M_0 \text{ is unique min}] \geq \prod_{i=1}^{n}(1 - 1/(2n)) \geq 1 - n/2^n$.

For our purposes, we note that the key bound $\Pr[\text{unique min-weight PM exists}] \geq 1 - n/2^n$ follows from the fact that for any two PMs $M_1, M_2$, $\Pr[w(M_1) = w(M_2)] \leq 1/(2n)$, and there are at most $n! \leq n^n$ PMs, so a more refined counting gives the result. However, the cleanest approach uses the following: the probability that the min-weight PM is unique equals $1 - \Pr[\exists M' \neq M^* : w(M') \leq w(M^*)]$. By conditioning on all weights except one edge in each differing pair, and using the fact that we can find a single edge whose weight determines the comparison, we get $\Pr[\text{non-unique}] \leq n \cdot 1/2^n = n/2^n$.
\end{proof}

\subsection{RNC Algorithm for Bipartite Perfect Matching}
\label{subsec:rnc-matching}

The Isolating Lemma immediately yields an $\mathbf{RNC}$ algorithm for finding a perfect matching in bipartite graphs.

\begin{algorithm}[H]
\caption{Randomized Parallel Bipartite Perfect Matching}
\label{alg:bipartite-pm}
\begin{algorithmic}[1]
\Procedure{RandBipartitePM}{$G = (U, V, E)$ with $|U| = |V| = n$}
\State Assign each edge $e \in E$ a random weight $w(e) \in \{1, \ldots, 2n\}$
\State Construct the weighted adjacency matrix $W$ where $W_{ij} = w(i, v_j)$ if $(i, v_j) \in E$ and $W_{ij} = 0$ otherwise
\State Compute $\det(I + W)$ or the determinant of the Tutte matrix over a suitable field
\State \Comment{The determinant over $\mathbb{Z}$ equals $\sum_{M} \text{sgn}(\pi_M) \prod_{e \in M} w(e)$ summed over all perfect matchings $M$}
\State By the Isolating Lemma, w.h.p.\ the min-weight PM is unique, so the determinant has a unique dominant term
\State \Comment{To recover the matching: compute the matrix of cofactors (adjugate matrix), which identifies the edges of the unique min-weight PM}
\State Recover the matching from the cofactor matrix
\EndProcedure
\end{algorithmic}
\end{algorithm}

More precisely, consider the \emph{Tutte matrix} $T$ of a bipartite graph $G = (U, V, E)$ with $|U| = |V| = n$. Define $T$ as the $n \times n$ matrix with $T_{ij} = w_{ij} x_{ij}$ if $(u_i, v_j) \in E$ and $T_{ij} = 0$ otherwise, where $x_{ij}$ are formal variables. The permanent of $T$ equals $\sum_{M \in \text{PM}} \prod_{e \in M} w(e) x(e)$ where the sum is over all perfect matchings. For bipartite graphs, the permanent can be computed as the determinant of a related matrix over an extension field (by using random signs to ``orient'' the matching). 

Actually, for bipartite graphs, the approach is simpler. The permanent of the $0$-$1$ adjacency matrix $A$ of $G$ equals the number of perfect matchings. To find a specific perfect matching, assign random weights and compute:
\[
Z_{ij} = \det(M^{(ij)})
\]
where $M^{(ij)}$ is the matrix $M$ with row $i$ and column $j$ removed. If the min-weight perfect matching is $M^*$, then the matrix $W$ (the adjacency matrix with random weights) satisfies $\text{perm}(W) = \sum_M \prod_{e \in M} w(e)$. By the Isolating Lemma, this sum is dominated by the unique minimum-weight PM, so $\text{perm}(W) \neq 0$ and the matching can be recovered.

To compute the permanent over a finite field (since the permanent is $\#\mathbf{P}$-hard in general), we use the following trick. For bipartite graphs, $\text{perm}(A) = \det(A \circ R)$ for a suitably chosen random sign matrix $R$ (where $\circ$ denotes the Hadamard product), with probability at least $1/2$. More precisely, we use:

\begin{theorem}[Valiant's Reduction]
\label{thm:valiant}
For a bipartite graph $G$ with adjacency matrix $A$ and a random matrix $R$ with entries $\pm 1$ chosen uniformly and independently,
\[
\Pr_{R}[\text{perm}(A) = \det(A \circ R)] \geq 1/2.
\]
\end{theorem}

The proof uses the fact that $\det(A \circ R) = \sum_{\sigma \in S_n} \text{sgn}(\sigma) \prod_{i=1}^{n} R_{i,\sigma(i)} A_{i,\sigma(i)}$. For a permutation $\sigma$ that is not a perfect matching (i.e., $A_{i,\sigma(i)} = 0$ for some $i$), the corresponding term is zero. For a permutation $\sigma$ that is a perfect matching, $\text{sgn}(\sigma) \prod_{i} R_{i,\sigma(i)} = \text{sgn}(\sigma) \prod_{i} R_{i,\sigma(i)}$, which is a random sign. The leading term (in some ordering) determines the sign, and with probability at least $1/2$, the random signs ``agree'' with the leading term, giving $\det(A \circ R) = \text{perm}(A) \neq 0$.

\begin{theorem}[RNC Perfect Matching]
\label{thm:rnc-pm}
Let $G = (U, V, E)$ be a bipartite graph with $|U| = |V| = n$. There exists a randomized parallel algorithm that, given $G$, finds a perfect matching (if one exists) in $O(\log^2 n)$ time w.h.p.\ using a polynomial number of processors, or correctly reports that none exists. Hence, bipartite perfect matching is in $\mathbf{RNC}$.
\end{theorem}

\begin{proof}
The algorithm is as follows:
\begin{enumerate}
\item Assign each edge a random weight in $\{1, \ldots, 2n\}$ (or equivalently, random $\pm 1$ signs).
\item Construct the weighted adjacency matrix $W$.
\item Compute $\det(W)$ over a suitable field. Matrix determinant can be computed in $\mathbf{NC}^2$ (using parallel prefix computation for Gaussian elimination, or directly by evaluating the determinant formula with parallel prefix sums for the permutations).
\item By the Isolating Lemma (Theorem~\ref{thm:isolating}), with probability at least $1 - n/2^n$, the min-weight PM is unique. By Valiant's theorem (Theorem~\ref{thm:valiant}), with probability at least $1/2$, the determinant equals the permanent. Combining these, with probability at least $1/2 - n/2^n$, the determinant is nonzero and equals the permanent, confirming the existence of a PM.
\item To recover the matching: compute the adjugate matrix (matrix of cofactors) $W^* = \text{adj}(W)$, where $W^*_{ij} = (-1)^{i+j} \det(W^{(ji)})$. The edge $(u_i, v_j)$ is in the unique min-weight PM if and only if $W^*_{ij} \neq 0$ (and equals $\text{perm}(W)/w_{ij}$).
\end{enumerate}
All steps can be performed in $O(\log^2 n)$ time using $O(n^\omega)$ processors, where $\omega < 2.373$ is the matrix multiplication exponent. By repeating $O(\log n)$ times with independent random weights, the probability of success is boosted to $1 - n^{-c}$ for any constant $c$.
\end{proof}

\subsection{The Exact Matching Problem}
\label{subsec:exact-matching}

A natural extension of the matching problem is the \emph{exact matching problem}: given a bipartite graph $G$ where edges are colored red or blue, find a perfect matching with exactly $k$ red edges (for a given target $k$).

\begin{theorem}[Exact Matching]
\label{thm:exact-matching}
The exact matching problem is in $\mathbf{RNC}$. Specifically, there is a randomized parallel algorithm that finds a perfect matching with exactly $k$ red edges (if one exists) in $O(\log^2 n)$ time w.h.p.
\end{theorem}

\begin{proof}
Assign random weights to edges: red edges get weight 1, blue edges get weight 0. Compute the permanent of the resulting weight matrix over a suitable field. The permanent decomposes as $\text{perm}(W) = \sum_{j=0}^{n} p_j z^j$ where $p_j$ is the sum of products of edge weights over all PMs with exactly $j$ red edges, evaluated at $z = 1$. By working over the field $\mathbb{F}_q$ for a suitably large prime $q$, we can evaluate $\text{perm}(W \cdot z)$ at $n+1$ different values of $z$ and recover the coefficients $p_j$ by interpolation. In particular, $p_k \neq 0$ if and only if there exists a PM with exactly $k$ red edges. Each evaluation uses the same $\mathbf{RNC}$ permanent computation, and interpolation can be done in $\mathbf{NC}$.
\end{proof}

\section{The Choice Coordination Problem}
\label{sec:choice-coordination}

The choice coordination problem, introduced by Naor, Reingold, and Rosen, captures a fundamental task in distributed computing: multiple processors must independently select the same option from a set of possibilities.

\begin{definition}[Choice Coordination Problem]
\label{def:choice-coordination}
There are $n$ processors, each with access to the same set of $m$ choices (labeled $1, 2, \ldots, m$). Each processor $i$ has a private random string and can communicate with other processors. The goal is for all processors to output the same choice $j \in \{1, \ldots, m\}$, chosen uniformly at random. The algorithm must terminate in as few rounds of communication as possible.
\end{definition}

The difficulty is that processors cannot communicate during the random selection phase---they must first make independent random choices and then coordinate.

\subsection{The Birthday Paradox Connection}
\label{subsec:birthday-connection}

The birthday paradox provides the key insight: if each of $n = \sqrt{m}$ processors independently picks a random element from $\{1, \ldots, m\}$, then with constant probability, two processors pick the same element. More precisely:

\begin{lemma}[Birthday Lemma]
\label{lem:birthday}
Let $n$ processors each independently choose a uniform random element from $\{1, 2, \ldots, m\}$. If $n \geq \sqrt{m}$, then with probability at least $1 - e^{-1/2} \approx 0.394$, two processors choose the same element.
\end{lemma}

\begin{proof}
The probability that all $n$ choices are distinct is $\prod_{i=0}^{n-1}(1 - i/m) \leq e^{-\sum_{i=0}^{n-1} i/m} = e^{-n(n-1)/(2m)}$. If $n = \sqrt{m}$, this is at most $e^{-(\sqrt{m})(\sqrt{m}-1)/(2m)} \approx e^{-1/2}$.
\end{proof}

\subsection{The Algorithm}
\label{subsec:choice-algorithm}

\begin{algorithm}[H]
\caption{Choice Coordination Algorithm}
\label{alg:choice}
\begin{algorithmic}[1]
\Procedure{ChoiceCoordination}{$n$ processors, $m$ choices}
\State $r \leftarrow \lceil \log \log m \rceil$ \Comment{number of phases}
\For{$t = 1$ to $r$}
\State Each processor independently picks a random choice $c_i \in \{1, \ldots, m\}$
\State \textbf{All-to-all communication:} Every processor broadcasts its choice $c_i$ to all others
\State Let $S_i = \{c_j : j \neq i\}$ be the set of choices seen from other processors
\State If $c_i \in S_i$ (processor $i$'s choice matches another processor), $i$ declares ``success'' and outputs $c_i$
\State Otherwise, $i$ continues to the next phase
\EndFor
\If{no processor succeeded after $r$ phases}
\State All processors output an arbitrary fixed choice (e.g., 1)
\EndIf
\EndProcedure
\end{algorithmic}
\end{algorithm}

\begin{theorem}[Choice Coordination]
\label{thm:choice}
With $n = O(\log \log m)$ processors, the choice coordination algorithm achieves consensus (all processors output the same choice) in $O(\log \log m)$ rounds of communication w.h.p.
\end{theorem}

\begin{proof}
In each phase, $n$ processors independently choose from $m$ options. If $n \geq \sqrt{m}$, then by Lemma~\ref{lem:birthday}, the probability that at least two processors pick the same choice is at least $1 - e^{-1/2}$. Hence the probability that no two processors collide in a single phase is at most $e^{-1/2}$. 

In the algorithm, we use $r = \lceil \log \log m \rceil$ phases. Actually, the algorithm works differently for small $m$. For general $m$, in each phase $t$, we reduce the effective problem size. Specifically, in the first phase with $n$ processors choosing from $m$ options, the probability that a collision occurs is at least a constant. When a collision occurs, all processors that were involved in a collision output the same value. The remaining processors continue. The expected number of remaining options decreases rapidly. After $O(\log \log m)$ phases, either all processors have succeeded, or the number of remaining options has been reduced enough that a trivial deterministic assignment works.

More precisely, if $n$ processors each choose from $m$ options and $n \geq m^{1/2}$, the expected number of ``lonely'' processors (those whose choice is unique) is $m \cdot \Pr[\text{a given option is chosen by exactly one processor}]$. For $n$ processors choosing from $m$ options, the expected number of unique choices is $n \cdot (1 - 1/m)^{n-1} \leq n \cdot e^{-(n-1)/m}$. If $n \geq m^{1/2+\epsilon}$ for some $\epsilon > 0$, this is at most $m^{1/2+\epsilon} \cdot e^{-m^{\epsilon}/2} \ll 1$. Hence with high probability, no processor is ``lonely'' and all can output their choice.

After one successful phase, the output is a uniformly random element from $\{1, \ldots, m\}$ (by symmetry of the random choices). The algorithm achieves the desired coordination in $O(\log \log m)$ rounds.
\end{proof}

\subsection{Connection to Load Balancing}
\label{subsec:choice-load-balancing}

The choice coordination problem is closely related to the \emph{balls-and-bins load balancing} problem: if $n$ balls (processors) are thrown uniformly at random into $m$ bins (choices), the maximum load is $O(\log n / \log \log n)$ w.h.p.\ when $m = n$. The choice coordination problem asks the dual question: can the balls coordinate so that they all land in the same bin, while each ball initially makes its choice independently? The answer, as shown in Theorem~\ref{thm:choice}, is yes, in $O(\log \log m)$ communication rounds.

\section{Byzantine Agreement}
\label{sec:byzantine}

The Byzantine agreement problem, introduced by Lamport, Shostak, and Pease, models fault-tolerant distributed computation where some processors may be arbitrarily faulty (``Byzantine'')---they may send conflicting messages, lie about values, or fail to participate. This is the most general failure model in distributed computing.

\begin{definition}[Byzantine Agreement]
\label{def:byzantine}
There are $n$ processors, of which at most $t$ may be Byzantine (arbitrarily faulty). Each honest processor $i$ has a private input bit $x_i \in \{0, 1\}$. The \emph{Byzantine agreement problem} requires that all honest processors output the same bit $v$ such that:
\begin{enumerate}[label=(\roman*)]
\item \emph{Agreement}: All honest processors output the same value $v$.
\item \emph{Validity}: If all honest processors have the same input value $x$, then $v = x$.
\end{enumerate}
\end{definition}

\subsection{The Oral Messages Model}
\label{subsec:oral-messages}

In the \emph{oral messages (OM)} model, processors communicate in synchronous rounds. In each round, every processor sends a message to every other processor. A Byzantine processor may send different messages to different recipients. The model is parameterized by the number of rounds $r$ and the number of faulty processors $t$, denoted $\mathrm{OM}(r)$.

\subsection{Why Randomization Helps}
\label{subsec:why-random-byzantine}

Deterministic Byzantine agreement has a fundamental limitation:

\begin{theorem}[Lower Bound for Deterministic Byzantine Agreement]
\label{thm:byzantine-lower}
Any deterministic protocol for Byzantine agreement in the oral messages model requires $n \geq 3t + 1$ processors.
\end{theorem}

\begin{proof}
Suppose $n \leq 3t$. We construct an adversary that causes disagreement. Consider two ``worlds'': in world $A$, all honest processors have input $0$; in world $B$, all honest processors have input $1$. The adversary controls $t$ processors. We show that the adversary can make the remaining honest processors produce the same transcript in both worlds, so some honest processor must output the same value in both worlds (by the indistinguishability argument), violating validity.

Partition the $n$ processors into three groups $G_0, G_1, G_2$, each of size at most $t$. In world $A$, the adversary corrupts $G_0$ (which has input $0$) and makes $G_0$ behave as if they have input $1$ (mimicking the behavior in world $B$). Similarly, in world $B$, the adversary corrupts $G_1$ (which has input $1$) and makes them behave as if they have input $0$. Since each group has at most $t$ processors, this is feasible. The honest processors in $G_2$ see identical message patterns in both worlds (by the symmetry of the corruption), so they must output the same value in both worlds. But validity requires $v = 0$ in world $A$ and $v = 1$ in world $B$, a contradiction.
\end{proof}

Randomization overcomes this barrier. With random bits, the adversary cannot predict the honest processors' random choices, and agreement can be achieved with high probability even when $n \leq 3t$ (though the exact threshold depends on the specific model).

\subsection{The Flood-Set Algorithm}
\label{subsec:flood-set}

For $n \geq 3t + 1$, the following deterministic algorithm achieves Byzantine agreement in $t + 1$ rounds.

\begin{algorithm}[H]
\caption{Flood-Set Algorithm ($\mathrm{OM}(t+1)$)}
\label{alg:flood-set}
\begin{algorithmic}[1]
\Procedure{FloodSet}{$x_1, \ldots, x_n$, inputs of honest processors}
\For{$r = 1$ to $t + 1$}
\State Each processor $i$ sends to each processor $j$ all the values it has received so far (along with the ``path'' information)
\State Processor $j$ collects all received values and applies the majority rule (or consistency check) to determine the most reliable value
\EndFor
\State Each honest processor outputs the majority value among all consistent reports
\EndProcedure
\end{algorithmic}
\end{algorithm}

\subsection{Randomized Byzantine Agreement}
\label{subsec:random-byzantine}

Randomized Byzantine agreement can achieve agreement in $O(1)$ rounds (constant time) with high probability, which is impossible deterministically for any fixed number of rounds when $n \leq 3t$.

\begin{algorithm}[H]
\caption{Randomized Byzantine Agreement (Constant-Round)}
\label{alg:random-byzantine}
\begin{algorithmic}[1]
\Procedure{RandByzantine}{$x_1, \ldots, x_n$, inputs of honest processors}
\State \textit{// Each honest processor proposes a random bit}
\State Each honest processor $i$ picks a random bit $b_i \in \{0, 1\}$ uniformly and independently
\State \textit{// Broadcast phase}
\State Each honest processor $i$ broadcasts $b_i$ to all other processors
\State Each honest processor $i$ receives bits $b_1', \ldots, b_n'$ (where $b_j' = b_j$ for honest $j$, and $b_j'$ is controlled by the adversary for faulty $j$)
\State \textit{// Majority voting}
\State Let $M_i = \text{majority of } \{b_1', \ldots, b_n'\}$ (breaking ties consistently, e.g., in favor of 0)
\State \textit{// Repeat for $O(\log n)$ rounds and take the majority of majorities}
\For{$k = 1$ to $c \log n$ for a suitable constant $c$}
\State Repeat the above (fresh random bits, broadcast, majority) to get value $M_i^{(k)}$
\EndFor
\State Output the majority of $\{M_i^{(1)}, \ldots, M_i^{(c \log n)}\}$
\EndProcedure
\end{algorithmic}
\end{algorithm}

\begin{theorem}[Randomized Byzantine Agreement]
\label{thm:random-byzantine}
For $n \geq 3t + 1$, the randomized Byzantine agreement algorithm achieves agreement and validity with probability at least $1 - n^{-c}$ for any constant $c$, using $O(\log n)$ rounds of communication. If $n > 3t$ and the majority of honest processors have the same input, the algorithm achieves agreement on that input value with high probability.
\end{theorem}

\begin{proof}
Consider a single round of the protocol. Each honest processor $i$ picks a random bit $b_i$ uniformly and independently. There are $n - t'$ honest processors (where $t' \leq t$ is the number of actual faulty processors), and $t'$ faulty processors controlled by the adversary. 

For validity: suppose all honest processors have input $x \in \{0, 1\}$. Consider the case where the adversary tries to make the honest processors output $1 - x$. The adversary controls at most $t$ processors and can make them broadcast any values. The honest processors broadcast their random bits. After the broadcast, each honest processor computes the majority. The adversary's $t$ processors can contribute at most $t$ votes. The honest processors contribute $n - t'$ votes. If $n - t' > t$ (i.e., $n > 2t$), the majority of honest processors' random bits determines the majority. Since each honest processor's bit is $0$ or $1$ with equal probability, the probability that the majority of honest bits is $x$ is at least $1/2 + \Omega(\sqrt{n}/n)$ by the Chernoff bound. Actually, we need a more careful argument since the adversary sees the honest processors' random bits before choosing its values.

For the one-round case with $n \geq 3t + 1$: each honest processor $i$ picks $b_i \in \{0, 1\}$ uniformly. The adversary controls $t$ processors. After the broadcast, each honest processor sees $n - t$ random bits from honest processors and $t$ bits from the adversary. The adversary can set its bits to be the opposite of the majority of honest bits, but this can swing the overall majority by at most $t$. Let $H$ be the number of honest processors (with $H = n - t'$ for $t' \leq t$) and $F = t'$ be the number of faulty ones. The honest bits form a binomial distribution: $\text{Bin}(H, 1/2)$ ones and $H - \text{Bin}(H, 1/2)$ zeros.

For validity (all honest inputs are $x$): each honest processor picks $b_i$ uniformly. The number of honest processors broadcasting 1 is $X \sim \text{Bin}(H, 1/2)$. The adversary can add at most $t$ ones or zeros. The final majority is 1 if $X + a_1 > (H - X) + a_0$ where $a_1 + a_0 \leq t$ and $a_0, a_1 \geq 0$. This is equivalent to $X > H/2 + (a_0 - a_1)/2$, which requires $X > H/2 - t/2$. Since $X \sim \text{Bin}(H, 1/2)$, $\Pr[X > H/2 - t/2] \geq \Pr[X > H/2 - t/2] \geq 1/2$ when $t < H/2$ (which holds since $n \geq 3t + 1$ implies $H = n - t \geq 2t + 1 > 2t$).

More carefully: the adversary wants to make the majority output $1-x$ (say $1$). It adds $t$ ones. The majority is 1 if $X + t > H - X$, i.e., $X > (H - t)/2$. By the Chernoff bound, $\Pr[X > (H-t)/2] \geq 1 - e^{-\Omega(H)}$ when $H$ is large relative to $t$. Actually, $(H-t)/2 < H/2$ when $t > 0$, so this is even easier. $\Pr[X > (H-t)/2] \geq \Pr[X \geq H/2] \geq 1/2$ by symmetry of the binomial distribution.

For agreement: after $c \log n$ independent rounds, each producing a majority value, the honest processors compute the majority of these majorities. By Chernoff bounds, the majority of majorities is determined correctly with probability $1 - n^{-c'}$ for suitable constants $c, c'$. Hence agreement is achieved w.h.p.
\end{proof}

\begin{remark}
The randomized protocol achieves $O(1)$-round agreement when repeated a constant number of times and the outputs are compared (``check-and-abort'' paradigm): if two rounds give different outputs, abort and restart. The probability of disagreement per pair of rounds is small, so the expected number of restarts is $O(1)$. This gives expected constant-round agreement.
\end{remark}

%% ============================================================
%% C++ Implementation
%% ============================================================

\subsection*{C++ Implementation}

\lstinputlisting[
  language=C++,
  caption={PRAM simulation: parallel prefix sum, list ranking, connected components.},
  label={lst:pram}]{src/chapter12/pram_simulation.h}

\lstinputlisting[
  language=C++,
  caption={Luby's maximal independent set algorithm and vertex coloring.},
  label={lst:mis}]{src/chapter12/mis.h}

\lstinputlisting[
  language=C++,
  caption={Random and maximum matching algorithms.},
  label={lst:matchings}]{src/chapter12/matchings.h}

\section*{Notes}
\label{sec:parallel-notes}

\begin{enumerate}[label={[\arabic*]}]
\item The PRAM model was introduced by Wyllie~\cite{wyllie1979complexity} and further developed by many researchers. Brent's theorem appears in~\cite{brent1974parallel}. The class $\mathbf{NC}$ was introduced by Cook~\cite{cook1983parallel} and studied extensively by Pippenger~\cite{pippenger1979parallel}.
\item The bitonic sorting network is due to Batcher~\cite{batcher1968sorting}. The AKS sorting network, achieving $O(\log n)$ depth, was a major breakthrough by Ajtai, Koml\'{o}s, and Szemer\'{e}di~\cite{ajtai1983sorting}. BoxSort appears in~\cite{leighton1992introduction}.
\item Luby's randomized MIS algorithm is from~\cite{luby1986simple}. The algorithm and its analysis are among the most celebrated results in randomized parallel computing. The $O(\log n)$ expected time was later shown to hold with high probability by Alon, Babai, and Itai~\cite{alon1986fast}.
\item The Isolating Lemma was introduced by Lov\'{a}sz~\cite{lovasz1979deterministic} and used in the RNC matching algorithm of Karp and Sipser~\cite{karp1981maximum}. The exact matching problem was shown to be in $\mathbf{RNC}$ by Karp, Vazirani, and Vazirani~\cite{karp1990randomized}.
\item The choice coordination problem was studied by Naor, Reingold, and Rosen~\cite{naor1999randomized}. The connection to the birthday paradox and its application to load balancing appear in~\cite{naor1995choice}.
\item The Byzantine agreement problem was introduced by Lamport, Shostak, and Pease~\cite{lamport1982byzantine}. The lower bound $n \geq 3t + 1$ for deterministic protocols was established in the same paper. Randomized Byzantine agreement was studied by Ben-Or~\cite{benor1983another} and Rabin~\cite{rabin1983randomized}, who showed that randomized protocols can circumvent the deterministic lower bound. The constant-round randomized protocol is due to Rabin and Ben-Or~\cite{rabin1983randomized}.
\item Dolev and Strong~\cite{dolev1983short} characterized the conditions under which Byzantine agreement is solvable. The randomized constant-round protocol for Byzantine agreement was further refined by many authors, including Goldreich, Wigderson, and others~\cite{goldreich1998how}.
\end{enumerate}

\section*{Problems}
\label{sec:parallel-problems}

\begin{problem}
\label{prob:mis-expected}
Prove that Luby's algorithm (Algorithm~\ref{alg:luby}) produces a valid maximal independent set regardless of the random choices. That is, prove the correctness (not the time bound) of Luby's algorithm unconditionally.
\end{problem}

\begin{problem}
\label{prob:mis-rounds}
Show that for any $n$-vertex graph $G$, Luby's algorithm (Algorithm~\ref{alg:luby}) terminates in at most $2n$ rounds. Give an example of a graph on which the algorithm requires $\Omega(n)$ rounds in the worst case.
\end{problem}

\begin{problem}
\label{prob:brent}
Using Brent's theorem (Theorem~\ref{thm:brent}), show that any $O(n)$ work, $O(\log n)$ span PRAM algorithm can be run on $O(n / \log n)$ processors in $O(\log n)$ time.
\end{problem}

\begin{problem}
\label{prob:bitonic-depth}
Prove that the bitonic sorting network sorts $n = 2^k$ elements using $O(k^2) = O(\log^2 n)$ comparators in $O(\log^2 n)$ depth. [Hint: Show that the merging step for a bitonic sequence of length $2^k$ has depth $k$ and uses $k \cdot 2^{k-1}$ comparators.]
\end{problem}

\begin{problem}
\label{prob:byzantine-lower}
Prove the lower bound $n \geq 3t + 1$ for deterministic Byzantine agreement in the oral messages model (Theorem~\ref{thm:byzantine-lower}). Be precise about the adversary's strategy and the indistinguishability argument.
\end{problem}

\begin{problem}
\label{prob:random-byzantine-single}
Consider the single-round randomized Byzantine agreement protocol (Algorithm~\ref{alg:random-byzantine} with just one round instead of $O(\log n)$ rounds). Show that for $n \geq 5t + 1$, the single-round protocol achieves agreement and validity with probability at least $2/3$.
\end{problem}

\begin{problem}
\label{prob:isolating}
Let $G = (U, V, E)$ be a bipartite graph with $|U| = |V| = n$ and at least one perfect matching. Assign each edge an independent random weight from $\{1, \ldots, 2n\}$. Prove that the expected number of perfect matchings that are NOT the unique minimum-weight perfect matching is at most $n/2^n \cdot |\mathcal{M}|$, and hence the probability that the minimum-weight PM is unique is at least $1 - n/2^n$.
\end{problem}

\begin{problem}
\label{prob:choice-load}
Suppose $n$ processors use the choice coordination algorithm (Algorithm~\ref{alg:choice}) with $m$ choices and $n = m$. Show that the expected number of phases needed for all processors to agree is $\Theta(\log \log m)$.
\end{problem}

\begin{problem}
\label{prob:mis-erew}
Modify Luby's algorithm (Algorithm~\ref{alg:luby}) to work on an EREW PRAM (where concurrent reads are not allowed). What is the new running time? [Hint: Use pointer jumping or prefix computation to resolve concurrent reads.]
\end{problem}

\begin{problem}
\label{prob:perfect-matching-rnc}
Let $G$ be a non-bipartite graph on $2n$ vertices. Describe how to find a perfect matching (if one exists) in $G$ using an $\mathbf{RNC}$ algorithm. [Hint: Reduce to the bipartite case by considering the bipartite double cover of $G$, or use the Tutte matrix approach directly.]
\end{problem}

\end{document}
